\section{The engine}
\label{sec:engine}

\subsection{Design principles}

The implementation follows three rules, each of which was adopted because
violating it caused a specific error during development.

\paragraph{One source of truth for parameters.} Every assumption is declared
once, in a schema carrying its symbol, unit, admissible range, grouping and
help text. Validation, the command line interface and the browser controls are
all generated from that schema. Nothing about the model is written down twice.
The alternative, a model that reads parameters and a user interface that
mirrors them, guarantees eventual divergence between the two.

\paragraph{Analytic and numerical layers coexist.} The closed forms of
\cref{sec:fold,sec:scaling} are implemented alongside the numerical solvers
that supersede them, and the test suite asserts agreement. The closed forms are
not documentation; they are the oracle against which the code is verified. This
is what caught the errors reported in \cref{sec:verification}.

\paragraph{Assumptions are outputs wherever possible.} The lumped constants
$f_s$, $S_m$ and $\FM$ remain as inputs to the algebraic layer, but the
first-principles layer computes them and the engine reports both side by side.
A model whose assumptions can be audited against its own detailed submodels is
much harder to fool oneself with.

\subsection{Module structure}

\begin{table}[t]
\centering
\caption{Modules and the sections of this paper they implement.}
\label{tab:modules}
\small
\begin{tabular}{@{}llp{6.6cm}@{}}
\toprule
Module & Sections & Responsibility \\
\midrule
\code{params.py}      & \ref{sec:notation}      & parameter schema, derived quantities, named presets \\
\code{sizing.py}      & \ref{sec:momentum} to \ref{sec:scaling} & momentum theory, closure, fold, sensitivities, fixed-point trace \\
\code{components.py}  & \ref{sec:blade}, \ref{sec:components}, \ref{sec:nearfield} & rotor blade model, beam-sized arms, motor and propeller mass, inertia, near-field corrections \\
\code{acoustics.py}   & \ref{sec:acoustics}     & source model, sound power, room field, A-weighting, blade passage \\
\code{electrical.py}  & \ref{sec:components}    & pack selection, $K_V$ sizing, currents, C-rate checks \\
\code{dynamics.py}    & \ref{sec:dynamics}      & state, forces, moments, allocation, hover linearisation \\
\code{control.py}     & \ref{sec:control}       & cascaded controller, authority, gain selection \\
\code{mission.py}     & \ref{sec:dynamics}, \ref{sec:control} & human trajectories, wind, reference governor with keep-out \\
\code{simulate.py}    & \ref{sec:dynamics}      & fixed-step integrator, telemetry, mission energy \\
\code{closure.py}     & \ref{sec:codesign}      & residual formulation, secant-accelerated battery loop, verification flight, margins \\
\code{numerics.py}    & \ref{sec:numerics}      & bounded mass-root search with explicit failure status, controllability rank \\
\code{optimize.py}    & \ref{sec:codesign}      & penalised objective, differential evolution, trade studies \\
\code{tether.py}      & \ref{sec:results}       & the closure with energy off-board \\
\code{sweep.py}       & \ref{sec:scaling}       & feasibility maps, exponent sweeps, endurance walls \\
\code{payloads.py}    & \ref{sec:results}       & display catalogue, readability, bill of materials \\
\bottomrule
\end{tabular}
\end{table}

\Cref{fig:modules} shows how the modules depend on one another, and
\cref{fig:pipeline} traces one evaluation of the coupled problem through them:
the order of operations that the rest of this paper derives piece by piece.

\begin{figure}[tbp]
\centering
% Import graph of the engine. An arrow A -> B means that module A imports B.
% Every module imports params.py, drawn as the base; the three interfaces at
% the top use most of the engine and their individual edges are omitted.
\begin{tikzpicture}[
    font=\footnotesize, y=1.05cm,
    mod/.style={draw, rounded corners=2pt, fill=cA!6, minimum width=1.9cm,
                minimum height=1.5em, font=\ttfamily\footnotesize},
    ifc/.style={mod, fill=cB!12},
    dep/.style={-{Stealth[length=1.8mm]}, cG, line width=0.55pt},
  ]
  % base
  \node[draw, fill=cG!12, minimum width=13.8cm, minimum height=1.5em,
        font=\ttfamily\footnotesize] (params) at (6.3,0)
        {params.py: schema, presets, derived quantities};
  % layer 1
  \node[mod] (numerics)  at (1.0,1)  {numerics};
  \node[mod] (acoustics) at (6.3,1)  {acoustics};
  \node[mod] (payloads)  at (8.6,1)  {payloads};
  \node[mod] (sizing)    at (11.6,1) {sizing};
  % layer 2
  \node[mod] (components) at (3.5,2) {components};
  % layer 3
  \node[mod] (electrical) at (0.8,3)  {electrical};
  \node[mod] (dynamics)   at (3.5,3)  {dynamics};
  \node[mod] (mission)    at (6.3,3)  {mission};
  \node[mod] (sweep)      at (10.0,3) {sweep};
  % layer 4
  \node[mod] (tether)     at (0.8,4) {tether};
  \node[mod] (control)    at (3.5,4) {control};
  % layer 5
  \node[mod] (simulate)   at (5.2,5) {simulate};
  % layer 6
  \node[mod] (closure)    at (2.8,6) {closure};
  \node[mod] (optimize)   at (7.6,6) {optimize};
  % interfaces
  \node[ifc] (api) at (2.6,7.7) {api};
  \node[ifc] (cli) at (6.3,7.7) {cli};
  \node[ifc, minimum width=2.6cm] (mk) at (10.4,7.7) {make\_results};
  \node (tool)  at (2.6,8.8)  {browser tool, \code{web/}};
  \node (term)  at (6.3,8.8)  {command line};
  \node (paper) at (10.4,8.8) {this paper};
  \draw[flow, cB] (api) -- (tool);
  \draw[flow, cB] (cli) -- (term);
  \draw[flow, cB] (mk) -- (paper);

  % engine edges: importer -> imported
  \draw[dep] (components.south west) -- (numerics.north east);
  \draw[dep] (dynamics.south west) -- (numerics.north);
  \draw[dep] (dynamics) -- (components);
  \draw[dep] (electrical.south east) -- (components.north west);
  \draw[dep] (mission.south west) -- (components.north east);
  \draw[dep] (sweep.south west) -- ($(components.east)+(0,0.12)$);
  \draw[dep] (sweep.south east) -- (sizing.north west);
  \draw[dep] (tether) -- (electrical);
  \draw[dep] (tether.south east) -- ($(components.north west)+(0.25,0)$);
  \draw[dep] (tether.west) -- ++(-0.55,0) |- (-0.35,0.52) -| (sizing.south);
  \draw[dep] (control) -- (dynamics);
  \draw[dep] (simulate.south west) -- (control.east);
  \draw[dep] (simulate.south) -- (dynamics.north east);
  \draw[dep] (simulate.south east) -- (mission.north);
  \draw[dep] (closure) -- (simulate);
  \draw[dep] ($(closure.south)+(-0.6,0)$) |- (components.west);
  \draw[dep] (closure.north) -- ++(0,0.25) -| ($(sizing.north)+(0.6,0)$);
  \draw[dep] (optimize) -- (simulate);
  \draw[dep] (optimize.south) |- ($(components.east)+(0,-0.12)$);

  % the interfaces use the whole engine
  \draw[decorate, decoration={brace, amplitude=5pt, mirror}, cB]
       ($(api.south west)+(-0.2,-0.12)$) -- ($(mk.south east)+(0.2,-0.12)$)
       node[midway, below=5pt, font=\scriptsize, fill=white]
       {the interfaces call across the whole engine};
\end{tikzpicture}

\caption{Import graph of the engine: an arrow from A to B means that module A
imports B. Every module also imports \code{params.py}, drawn as the base. The
layers read from the bottom up as the layers of the paper: the algebraic and
component models, the dynamics and control, the simulation, and the two
modules that close the coupled problem, \code{closure.py} and
\code{optimize.py}. \code{acoustics.py} and \code{payloads.py} import nothing
but the parameters. The three interfaces at the top call across the whole
engine and their individual edges are omitted.}
\label{fig:modules}
\end{figure}

\begin{figure}[tbp]
\centering
% How the coupled problem is solved: a bounded root search sizes the hover
% design, the battery loop wraps an inner algebraic fixed point and one
% simulation per update, a verification flight closes it, and the results
% flow into this paper.
\begin{tikzpicture}[
    font=\footnotesize,
    box/.style={block, text width=2.8cm, minimum height=3.3em, inner sep=3pt},
    stp/.style={box, fill=cA!7},
    sim/.style={box, fill=cB!12},
    res/.style={box, fill=cC!10},
    lab/.style={font=\scriptsize},
  ]
  % top: hover sizing
  \node[box] (p)    at (0,0)    {parameter set $p$,\\ \code{params.py}};
  \node[stp] (root) at (3.9,0)  {component closure\\ $\mathcal{F}(m)=0$ by a\\ bounded root search};
  \node[box] (hov)  at (7.8,0)  {hover design\\ $m,\ \mbat^{(0)}$};
  \draw[flow] (p) -- (root);
  \draw[flow] (root) -- (hov);

  % middle: the battery loop
  \node[stp] (inner) at (0,-2.6)   {inner map $\mathcal{G}(\mu)$:\\
       $m=\mfix+S(m)+\mu$\\ by contraction};
  \node[stp] (veh)   at (3.9,-2.6) {vehicle: $J$, $\mat M$,\\ $k_T,k_Q,k_P$, trim};
  \node[sim] (fly)   at (7.8,-2.6) {simulate a window:\\ RK4, $h=\SI{4}{\milli\second}$,\\ controller, governor};
  \node[stp] (req)   at (7.8,-5.1) {requirement\\ $\max\{E/\eb(1{-}\varsigma),$\\ $P_{\mathrm{peak}}/p_b\}$};
  \node[stp] (sec)   at (3.9,-5.1) {secant step\\ on $\mathcal{R}_2$, guarded,\\ \labelcref{eq:secant}};
  \node[box] (test)  at (0,-5.1)   {$|\mathcal{R}_2|<\mathrm{tol}$\,?};
  \draw[flow] (hov.south) -- ++(0,-0.55)
       -| node[lab, pos=0.3, below] {$\mu_0=\mbat^{(0)}$} (inner.north);
  \draw[flow] (inner) -- node[lab, above] {$m$} (veh);
  \draw[flow] (veh) -- (fly);
  \draw[flow] (fly) -- node[lab, right] {$\bar P,\ P_{\mathrm{peak}}$} (req);
  \draw[flow] (req) -- node[lab, above] {$\mathcal{H}(\mu)$} (sec);
  \draw[flow] (sec) -- node[lab, above] {$\mathcal{R}_2$} (test);
  \draw[flow] (test.north) -- node[lab, left] {no} (inner.south);

  \begin{scope}[on background layer]
    \node[draw=cA, dashed, rounded corners=3pt, inner sep=7pt,
          fit=(inner)(veh)(fly)(req)(sec)(test)] (loop) {};
  \end{scope}
  \node[lab, cA, anchor=north east] at (loop.south east)
       {battery loop, \code{closure.py}: one simulation per update};

  % bottom: verification and output
  \node[sim] (ver) at (0,-7.9)    {verification flight\\ at the rounded battery};
  \node[res] (con) at (3.9,-7.9)  {margins and path\\ constraints \labelcref{eq:cd4}};
  \node[res] (res) at (7.8,-7.9)  {\code{make\_results.py}:\\ results, macros,\\ tables, figure data};
  \node[res] (pap) at (11.5,-7.9) {this paper};
  \node[res] (api) at (11.5,-2.6) {API and\\ browser tool};
  \draw[flow] (test.south) -- node[lab, left] {yes} (ver.north);
  \draw[flow] (ver) -- (con);
  \draw[flow] (con) -- (res);
  \draw[flow] (res) -- (pap);
  \draw[flow, cB] (p.north) -- ++(0,0.4) -| (api.north);
  \draw[flow, cB, dashed] (ver.west) -- ++(-0.35,0) |- ($(ver.south)+(0,-0.45)$)
       node[lab, pos=0.75, below] {round up, fly again} -- (ver.south);
\end{tikzpicture}

\caption{One evaluation of the coupled problem as the engine carries it out.
A bounded root search on the component residual \cref{eq:fpclosure} gives the
hover design, which seeds the battery loop. Each update of the loop solves the
inner mass map by contraction (\cref{prop:inner}), rebuilds the vehicle, flies
one simulation window and turns its mean and peak power into a battery
requirement by \cref{eq:batreq}; a guarded secant step then moves the battery.
When the residual is within tolerance the battery is rounded up and flown
again, and the loop ends only on a verification flight that needs no more
battery than it carries. The results pass through \code{make\_results.py} into
this paper, and the same engine serves the API and the browser tool.}
\label{fig:pipeline}
\end{figure}

\subsection{Interfaces}

Three front ends share the engine: a command line for scripted studies, an HTTP
API, and a browser tool that generates its controls from the schema and
recomputes on every parameter change. The tool exists because the questions
that matter here, ``what does this assumption cost'', are answered by moving one
number and watching six others, which is a different activity from running a
batch job.

The browser tool renders the assembly in three dimensions from the same solved
dimensions used by the closure: arm length and wall thickness from
\cref{eq:armlength,eq:wall}, rotor diameter and chord from
\cref{sec:blade}, motor body size from \cref{eq:motormass}, and the individual
battery cells from the pack selection of \cref{sec:components}. The check that
this is a drawing of the machine rather than a picture beside it is that the
drawn parts sum to the closure's gross mass.

\subsection{Reproducing the results}

Every number in \cref{sec:results} is produced by one script,
\code{paper/make\_results.py}, which writes the numbers together with the
configuration that produced them (preset, payload, window, step, settling
interval, seed, governor limits, solver tolerances and optimiser settings) to
\code{paper/results/results.json}, and emits the table rows and macros that the
\LaTeX\ sources include. No figure in the design study is typed by hand.
\code{paper/build.py} then compiles this document.
